\documentclass[11pt,a4paper]{article}
\usepackage[T1]{fontenc}
\usepackage[utf8]{inputenc}
\usepackage[margin=2.05cm]{geometry}
\usepackage[expansion=false]{microtype}
\usepackage{enumitem,amsmath}
\setlist[itemize]{leftmargin=1.1em,itemsep=1pt,topsep=2pt}
\setlist[enumerate]{leftmargin=1.5em,itemsep=1pt,topsep=2pt}
\usepackage[hidelinks]{hyperref}
\pagestyle{empty}
\setlength{\parskip}{2pt}
\title{\vspace{-1.6cm}\large\bfseries A pre-registered \emph{in silico} falsification programme\\
on a redox-switch model of parvalbumin-interneuron damage}
\author{\normalsize Nikita Filippov \quad\textperiodcentered\quad \href{https://github.com/stockman2/SchizoStudy-v5}{https://github.com/stockman2/SchizoStudy-v5} \label{row:target}}
\date{\normalsize 23 October 2026}
\begin{document}
\maketitle
\vspace{-0.9cm}

\noindent\textbf{What was done.} We built a mechanistic ODE model of oxidative damage to
parvalbumin interneurons --- a bistable peroxiredoxin/sulfiredoxin redox switch, a two-stage
inducible defence ($\tau_m=12$\,h, $\tau_s=63$\,h), a Phillips--Robinson sleep--wake switch, and a
Wilson--Cowan cortex closing the loop through disinhibition --- and then spent most of the effort
trying to break it. Every hypothesis was written as a numbered rule with its parts, and
\emph{every threshold was derived from expectations and variances before the computation ran}.
Verdicts were recorded in an append-only ledger and were never revised afterwards; when a defect
was found in a rule, the defect was recorded \emph{beside} the verdict and the question was re-asked
as a new rule on a fresh noise seed.

Of \textbf{127 pre-registered parts across 26 rules: 66 passed, 47 failed, 13 were unfalsifiable,
and one fell outside its own analytic bracket.} More than a third of our own predictions were
wrong. That ratio is the main result we wish to report, and the failures below were more
informative than the successes.

\vspace{2pt}
\noindent\textbf{What survived.}
\begin{enumerate}
\item \emph{Episode duration is set by defence memory, not by repair capacity.} Within a parameter
set, 99.8\,\% of the variance of episode duration is explained by the history of oxidative load and
0.2\,\% by repair; the decomposition is exact by construction and the shares sum to unity to
$10^{-16}$. The dependence on defence memory is affine and does not saturate: a saturating form
with a shared offset is rejected, although our quantitative out-of-fit prediction missed its own
$2\,\mathrm{SE}$ band. Consequence: passively strengthening repair cannot reach the clinical
duration target --- it misses by a factor of 2 to 10.
\item \emph{Slow defence memory generates a pre-episode precursor.} Removing it lowers the
precursor in 107 of 174 sets, $p=0.0015$. Slow load noise is not needed for this; its two
independent justifications were withdrawn.
\item \emph{The population lives in three phases, not two.} Between the healthy core and the
collapsed core there is a band of cells that are \emph{more} vulnerable than their collapsed
neighbours and nevertheless intact --- roughly 8.5\,\% of the population, seventeen standard
errors from zero. It collapses to \emph{exactly} zero when the slow defence is frozen, so it is a
dynamical effect rather than a parametric one, and priming widens it without causing it. This is
hormesis at the population level and it is a prediction of the model, not an artefact of counting.
\item \emph{The rhythm breaks through a devil's staircase}, with rational rotation numbers,
alternating cycle lengths on the fractional steps, and a cascade that adds one rather than
doubling.
\item \emph{Cell count enters only as effective spread of resilience}, not as numerical resolution;
the $J\to\infty$ limit exists and was measured, with corrections converging as $1/J$.
\item \emph{A capacity-dependent feedback loop converts frequent short episodes into rare long
ones.} With per-cell capacity loss while damaged and slower recovery afterwards, mean episode
duration rises from 13.1 to 25.6 days while the onset rate \emph{falls} from 0.042 to 0.033 per
at-risk set-day; episodicity is retained at 0.75 of the open model; the broken state is held after
load is released, monotonically in release depth ($\rho=+1.000$, $p=0.0083$, the attainable
minimum); and progression exceeds a capacity-matched static control by $+0.0231$ against a minimum
detectable effect of $0.0096$ --- replicated as $+0.0223$ and $+0.0211$ on two further independent
noise seeds.
\end{enumerate}

\vspace{2pt}
\noindent\textbf{What failed.}
\begin{itemize}
\item \emph{A load-dependent feedback loop of the Grace type is refuted as a model of the
disease.} Driving oxidative load with the damaged fraction does produce irreversibility, and it is
not a mere gain knob: it latches states that no static load reaches at all. But irreversibility
switches on at loop strength 1.11 while episodes extinguish at 0.18 --- a sixfold gap that is
structural rather than marginal, and robust to a gentler breaking protocol. There is no window in
which chronicity and episodicity coexist. The reason is now clear: load is a per-set quantity, so a
loop acting through it moves the whole population together.
\item \emph{The capacity loop changes the level and the duration of damage but not the shape of
risk in time.} Its hazard trend exceeds a capacity-matched static arm by only $1.94$ against a
threshold of $2.23$ computed exactly for that comparison.
\item \emph{Bursty load does not beat smooth load}, contrary to the hypothesis that rare strong
excursions dominate; and the measured quantity that first appeared to support it turned out to
saturate.
\item Several precursor-slope predictions were \emph{unfalsifiable} at the grid resolution we had
chosen: the minimum detectable effect was larger than the threshold we had written down.
\end{itemize}

\vspace{2pt}
\noindent\textbf{What the failures taught (the transferable part).} The single most useful finding
is negative and methodological:

\emph{In a model with absorbing states and a heterogeneous population, no absolute null exists for
any accumulation measure.} The damaged fraction rises on its own, because a bistable switch does
not release a collapsed set until load is lowered; and the hazard of new episodes declines on its
own, because the most vulnerable cells are absorbed first and the survivors are less vulnerable.
Two of our own controls asserted that these baselines were zero. Both were false, and --- this is
the instructive part --- both \emph{passed at small sample size and failed as data accumulated},
because the baseline is fixed while the detection threshold shrinks as $1/\sqrt{n}$. A control whose
null is false is passed by lack of power. Every claim must therefore be a matched difference
between arms brought to the same damage level, never a test against zero.

Nine further statistical traps were catalogued, each of which had once produced a false conclusion
in our own hands and each of which is now closed by a check in the code. Four are worth naming:
a variance-ratio statistic used in four rules was shown analytically to equal $1+\mathrm{CV}^2$ of
the underlying rates and was declared unfit; an equivalence claim whose tolerance exceeds half the
range of its own curve cannot be falsified; a threshold must be \emph{attainable}, since $p\le0.05$
on an exact permutation test over three points has a floor of $1/3!=0.167$, and we once recorded
``failed'' for a perfectly monotone result; and a comparison point chosen as the nearest element of
a grid that does not bracket the target answers a question about the grid's edge.

\vspace{2pt}
\noindent\textbf{Scope, stated as a result rather than a caveat.} The model has no absorption-free
stationary state, so its domain of validity is windows short compared with the accumulation time.
It is therefore a model of the \emph{generation and termination of an episode} under oxidative
load, not of the course of the illness. We did not attempt the global picture and do not claim it.
The capacity loop is the only mechanism we tested that makes progression appear at all, and it does
so by lengthening episodes rather than by making them more frequent.

\vspace{2pt}
\noindent\textbf{Reproducibility.} A fused compiled kernel giving a 3.3-fold wall-clock speed-up was
proved identical to the reference integrator at two levels: bitwise over 67 output fields on nine
parameter pins (worst deviation $9.5\times10^{-15}$ against a tolerance of $10^{-12}$), and over the
entire set of verdicts. Block results are cached under a key that includes the hash of the model
source, the kernel included --- a discipline added after a kernel defect was found that no pin
touched. Step-size convergence separates two classes of quantity: continuous means converge to
within 0.5\,\% by $M=800$, whereas counting quantities do not, with a Richardson limit 2.0\,\% above
the value at $M=1600$; counting quantities are therefore not tested against thresholds finer than
that. Whole project could be found in the \hyperref[row:target]{github}.

\vspace{2pt}
\noindent\textbf{What we would ask next.} The load loop failed because load is a per-set quantity,
so a loop acting through it moves the whole population together; a structure in which regions couple
through narrow bottlenecks would break that homogeneity, and is the natural next test. The capacity
loop lengthens episodes but leaves the shape of risk unchanged, which suggests that progression on
the scale of the illness needs a component that does \emph{not} fully return. We have deliberately
not added it: a mechanism introduced to produce the result it was designed to produce is not a test,
and the rule that would falsify it has to be written first.

\end{document}
